\chapter{Les fonctions de transitions utilis\'ees}
\mysetref{the-rules}

\label{rule-xor9}%
La r\`egle~\arule{Xor9}:\\
$\Center= \Center \xor \north \xor \south \xor \east \xor \west \xor\\
 \Neast \xor \Nwest \xor \Swest \xor \Seast$

\mycodefull{c-qmean}%
{}%
{Le c{\oe}ur de la r\`egle~\arule{Qmean}}%
{\label{rule-qmean}%
  Nous avons exp\'eriment\'e plusieurs param\`etres statiques\footnote{%
    \mygetref[phrase]{static-param}} %
  pour cette r\`egle qui nous a servi initialement \`a tester le moteur \`a
  voisinage sp\'ecial d\'ecrit en~\mygetfig{func-7}. %
  Les param\`etres utilis\'es dans les illustrations de
  la~\mygetref{obj-visualisation} sont:~%
  mask = 3; match = 0; delta = -1. %
  Ces param\`etres sont parmi les premiers que nous ayons essay\'es, et
  donne un des comportements les plus riches\footnote{%
    Intuitivement, en remarquant cependant que
    ce para\-m\'e\-trage est ins\-pir\'e par la recherche d'une opposition
    \bialt{cons\-truc\-tion}{des\-truc\-tion},
    \bialt{\'egal\-isa\-tion}{dif\-f\-\'er\-en\-cia\-tion}, etc.\  typique d'un r\`egle
    telle~\arule{Life}, et \'evocatrice de ph\'enom\`enes de fronti\`eres.} %
  que nous ayons trouv\'es pour cette famille\footnote{%
    La famille est ici comprise comme l'espace de r\`egles engendr\'e par
    un syst\`eme de param\`etres statique pour un algorithme de g\'en\'eration de
    table de transition.} %
  de r\`egles.}

\mycodefull{c-anneal}%
{}%
{La r\`egle~\arule{Anneal} est impl\'ement\'ee via la r\`egle~\arule{Sum9}}%
{\label{rule-anneal}%
  Le source de la fonction \asoft{C} impl\'ementant la fonction de
  transition de la r\`egle~\arule{Sum9}, qui re\c{c}oit en argument statique
  une table secondaire d\'ecrivant une r\`egle additive. %
  \mygetref[phrase]{additive-rules}}%

\mycodefull{tcl-anneal}%
{}%
{L'interface de binding \asoft{Tcl} pour~\arule{Anneal}}%
{\arule{Anneal} est d\'ecrite sous les noms:~%
    \arule{M46789} dans~\cite{Vichniac86}, \\%
    \arule{Lava Lamp}~(une version~3D) dans~\cite{Pickover92b}, %
    \arule{Anneal} dans~\cite{Toffoli87}. %
    Il s'agit d'une modification de la r\`egle
    dite~\arule{Vote~Majoritaire}~\cite{Toffoli87}, qui
    \aquote{recuit} la dynamique du vote majoritaire par \'echange des
    bits~5 et~6.}

\mycodefull{c-tube-worm}%
{}%
{La r\`egle~\arule{Tube-Worm}}%
{\label{rule-tube-worm}%
  Le source de la fonction \asoft{C} impl\'ementant la fonction de
  transition de la r\`egle~\arule{Tube-Worm},
  d'apr\`es~\cite{Toffoli87}.}

%%
%% Local Variables:
%% mode: latex
%% mode: auto-fill
%% iso-accents-minor-mode: nil
%% TeX-master: "top"
%% TeX-command-default: "mytex"
%% Local IspellParsing: "latex-mode ~latin1"
%% Local IspellPersDict: "~/.ispell_lisp"
%% LocalWords: "the-rules Xor c-qmean ur Qmean static-param func mask m\'e Sum"
%% LocalWords: "obj-visualisation trage c-anneal Anneal additive-rules binding"
%% LocalWords: "tcl-anneal Tcl Lamp c-tube-worm Tube-Worm"
%% End:
